%-------------------------------------------------------------------
% Appendix: RSA-based numerical illustration of d_cog
% This appendix can be inserted into topology_grounding.tex or
% grounding_without_reality.tex to address the operationalisability
% of d_cog.
%-------------------------------------------------------------------

\appendix
\section{Numerical Illustration of $d_{\mathrm{cog}}$ via RSA}
\label{app:dcog-numerical}

We demonstrate that $d_{\mathrm{cog}}$ is numerically tractable via
Representational Similarity Analysis (RSA)~\cite{kriegeskorte2008},
addressing the concern that the Gromov--Hausdorff formulation is
mathematically defined but empirically inaccessible.

\subsection*{Setup}

Let $\Sigma = \{\textsc{flower}, \textsc{food}, \textsc{same},
\textsc{different}, \textsc{zero}\}$. We consider three agents:

\begin{itemize}[noitemsep]
  \item $\mathcal{A}_{\mathrm{bee}}$: honeybee, UV- and
    olfactory-dominant sensory space;
  \item $\mathcal{A}_{\mathrm{hum}}$: human, colour-opponent and
    proprioceptive-dominant;
  \item $\mathcal{A}_{\mathrm{llm}}$: a multimodal language model,
    text-cooccurrence-dominant with weak sensory channels.
\end{itemize}

Each agent's grounding map $G: \Sigma \to \mathcal{K}(\mathcal{S})$ is
approximated by a centroid vector in a six-dimensional sensory channel
space (UV, olfactory, colour-opponent, tactile, proprioceptive,
text-cooccurrence). The RSA matrix $\mathrm{RSA}_{\mathcal{A}}$ is the
cosine-similarity matrix over these centroid vectors, and the induced
metric on $\Sigma$ is $d_{\mathcal{A}}(\sigma_i, \sigma_j) =
\sqrt{1 - \mathrm{RSA}_{\mathcal{A}}(i,j)}$ (angular distance).

\subsection*{Results}

Table~\ref{tab:dist} shows the induced pairwise distances on $\Sigma$.

\begin{table}[ht]
\centering
\small
\caption{RSA-induced distances on $\Sigma$ for three agents.
Entries are $d_{\mathcal{A}}(\sigma_i,\sigma_j)$ (angular distance,
range $[0,1]$).}
\label{tab:dist}
\begin{tabular}{l|ccccc}
\hline
\textbf{Bee} & flower & food & same & diff. & zero \\
\hline
flower    & 0.000 & 0.229 & 0.505 & 0.505 & 0.244 \\
same      & 0.505 & 0.388 & 0.000 & \textbf{0.000} & 0.310 \\
\hline
\textbf{Human} & & & & & \\
\hline
flower    & 0.000 & 0.344 & 0.745 & 0.727 & 0.831 \\
same      & 0.745 & 0.705 & 0.000 & \textbf{0.025} & 0.160 \\
\hline
\textbf{LLM} & & & & & \\
\hline
flower    & 0.000 & 0.131 & 0.297 & 0.297 & 0.343 \\
same      & 0.297 & 0.337 & 0.000 & \textbf{0.000} & 0.112 \\
\hline
\end{tabular}
\end{table}

Approximating $d_{\mathrm{GH}}$ by exhaustive minimisation over all
$5! = 120$ bijections $f: \Sigma \to \Sigma$:

\[
d_{\mathrm{cog}}(\mathcal{A}_{\mathrm{bee}},\,\mathcal{A}_{\mathrm{hum}})
  = 0.208, \quad
d_{\mathrm{cog}}(\mathcal{A}_{\mathrm{hum}},\,\mathcal{A}_{\mathrm{llm}})
  = 0.238, \quad
d_{\mathrm{cog}}(\mathcal{A}_{\mathrm{bee}},\,\mathcal{A}_{\mathrm{llm}})
  = 0.099.
\]

\subsection*{Observations}

Three structural findings are noteworthy.

\textbf{(1) \textsc{same}/\textsc{different} collapse.}
Both the bee and the LLM assign distance $0.000$ to the pair
$(\textsc{same}, \textsc{different})$, but for structurally distinct
reasons. In the bee, the olfactory--UV space provides no dimension along
which the relational contrast is grounded sensorimotorly; behavioural
evidence for the bee's same/different discrimination
\cite{giurfa2001} therefore reflects a functional capacity that is not
accompanied by topological separation in the grounding-induced metric.
In the LLM, the collapse arises from the absence of a closed
action-consequence loop: without resistance, the two symbols are
distinguished only by text-cooccurrence patterns, which are nearly
symmetric. The human exhibits a small but nonzero separation
($d = 0.025$), reflecting proprioceptive grounding of relational
matching.

\textbf{(2) Human--LLM gap is the largest.}
$d_{\mathrm{cog}}(\mathcal{A}_{\mathrm{hum}},
\mathcal{A}_{\mathrm{llm}}) = 0.238 >
d_{\mathrm{cog}}(\mathcal{A}_{\mathrm{bee}},
\mathcal{A}_{\mathrm{hum}}) = 0.208$.
The LLM's grounding topology is structurally more distant from the
human than the bee is, despite the LLM having nominal access to visual
modalities. This supports the claim of Section~4.2 that adding sensory
channels does not automatically reduce $d_{\mathrm{cog}}$: the
cross-modal integration architecture matters.

\textbf{(3) Operationalisability.}
The computation above uses only (i) a choice of sensory channel
dimensions, (ii) a method for estimating grounding centroids (here,
assumed; in practice, neural activation patterns via RSA), and (iii)
angular distance and exhaustive bijection search. For larger
$|\Sigma|$, the bijection search can be replaced by approximation
algorithms for GH distance \cite{memoli2007}. The calculation is
therefore tractable in principle with standard neuroscientific
measurements.

\paragraph{Limitations.}
The channel vectors used here are illustrative assumptions, not
empirical measurements. The point of this appendix is to demonstrate
that $d_{\mathrm{cog}}$ admits numerical instantiation, not to make
empirical claims about these specific agents. Empirical instantiation
via RSA on neural or model activation data is the subject of OP1.

\paragraph{Note on higher-dimensional components.}
The five-concept inventory used here is sufficient to illustrate the
$d_{\mathrm{GH}}$ and $d_{\mathrm{bn}}^{(0)}$ components of the
persistence signature
(Definition~\ref{def:dcog-signature}), but cannot support
$d_{\mathrm{bn}}^{(k)}$ for $k \geq 1$: the Vietoris--Rips filtration
of a five-point metric space produces at most trivial $H_1$ diagrams.
Appendix~\ref{app:dcog-h1} extends the construction to a sixteen-concept
inventory at which $H_1$ features emerge, and characterises the
detection threshold quantitatively.
